\documentclass[a4paper]{article}

\usepackage[utf8]{inputenc}

\usepackage{amssymb, amsmath, amsthm}
\usepackage{bbold}
\usepackage{graphicx}
\usepackage{enumitem}
\usepackage{amsfonts}
\usepackage{mathtools}
\usepackage{orcidlink}
\usepackage{colonequals}
\usepackage{lmodern}
\usepackage{microtype}
\usepackage{mathrsfs}

\usepackage{algorithm2e}

% Bibliography settings
%---------------------------
\usepackage[style=ext-numeric,
            maxbibnames=99,  % Never write 'et al.' in a bibliography
            maxcitenames=99,  % Never write 'et al.' in a citation
            giveninits=true, % Abbreviate first names
            sortcites=true,  % Sort citations when there are more than one
            backend=biber]{biblatex}
\addbibresource{dune-gfe-manual.bib}

% Transform certain titles to "sentence case".  Needs an ext- style.
% See https://tex.stackexchange.com/questions/22980/sentence-case-for-titles-in-biblatex
\DeclareFieldFormat
[article,inbook,incollection,inproceedings,patent,thesis,unpublished]
{titlecase:title}{\MakeSentenceCase*{#1}}

% Put a thin space between two consecutive initials
\renewrobustcmd*{\bibinitdelim}{\,}

% Settings for the listings package
%--------------------------------------

\usepackage{listings}
\lstset{language=[11]{c++},
         basicstyle=\small\ttfamily,
         commentstyle=\textit,
         columns=flexible,keepspaces=true,
         showstringspaces=false,      % no special string spaces
         escapeinside={/*@}{@*/}
        }

% How to include ranges of an external source code file
\lstset{rangeprefix=//\ \{\ ,% curly left brace plus space, all in a C++-style comment
        rangesuffix=\ \},% space plus curly right brace
        numberstyle=\footnotesize,  % font size for numbers
        includerangemarker=false}  % Do not show the range markers

\definecolor{interfacecolor}{rgb}{0.95,0.95,1}
\definecolor{interfaceclasscolor}{rgb}{0.98,0.81,0.81}
\definecolor{cmakecommandcolor}{rgb}{0.2,0.8,0.8}
\lstdefinestyle{Example}{}
\lstdefinestyle{Interface}{backgroundcolor=\color{interfacecolor},frame=single}
\lstdefinestyle{InterfaceClass}{backgroundcolor=\color{interfaceclasscolor},frame=single}
\lstdefinestyle{CMakeCommand}{backgroundcolor=\color{cmakecommandcolor},frame=single}

\newcommand{\cpp}[1]{\lstinline[basicstyle=\ttfamily]!#1!}

\lstnewenvironment{cppinterface}[1][]{\lstset{style=Interface}}{}
\lstnewenvironment{cppinterfaceclass}[1][]{\lstset{style=InterfaceClass}}{}

% For typesetting Dune module names
\newcommand{\dunemodule}[1]{\texttt{#1}}

%  Custom environments
%---------------------------------

\newtheorem{definition}{Definition}
\newtheorem{lemma}{Lemma}
\newtheorem{remark}{Remark}



\usepackage{todonotes}
\newcommand{\todosander}[1]{\todo[inline,color=white,author=OS,caption={}]{#1}}

\usepackage{float}

\usepackage{tikz}
\usetikzlibrary{arrows}

\usepackage{multirow}
\usepackage{changes}
\usepackage{hyperref}

\newcommand{\R}{\mathbb R}
\newcommand{\N}{\mathbb N}

\DeclarePairedDelimiter{\abs}{\lvert}{\rvert}
\DeclarePairedDelimiter{\norm}{\lVert}{\rVert}


% Identity matrix
\newcommand{\identity}{\mathbb{1}}
\newcommand{\Tref}{{T_\textnormal{ref}}}


\DeclareMathOperator{\dist}{dist}
\DeclareMathOperator{\SOdrei}{SO(3)}
\DeclareMathOperator*{\argmin}{arg\,min}

%%  All graphics files may be in this subdirectory
\graphicspath{{gfx/}}



\title{The dune-gfe manual}
\author{Oliver Sander}

\begin{document}

\maketitle

\begin{abstract}
 \dunemodule{dune-gfe} is a module in the \textsc{Dune} ecosystem that implements
 various geometric finite element (GFE) methods.
 This document collects miscellaneous technical information about the implementation.
\end{abstract}

\section{Geometric finite elements}

Let $\Omega$ and $M$ be smooth manifolds.  Geometric finite elements provide approximations for maps
$v : \Omega \to M$ by suitably generalizing piecewise polynomials.  For this,
$\Omega$ is discretized by a grid $\mathcal{T}$ with elements $T$.
These elements can be simplices, (hyper-)cubes, or more exotic shapes.
For details about the construction of geometric finite elements see
\cite{sander:2010,sander:2012,sander:2015,grohs_hardering_sander_sprecher:2019,hardering_sander:2020}.

\subsection{Representation of target manifolds}
\label{sec:representation_of_manifolds}

Currently, the module assumes that $M$ is embedded into $\R^n$.  In other words,
elements of $M$ are stored as points in $\R^n$. In the code, these coordinates
are typically called \emph{global coordinates}, and they are available via the
\cpp{globalCoordinates()} method. For tangent vectors there are two representations:
One can either store tangent vectors as vectors in $\R^n$ as well.  In this case
they are referred to as \emph{embedded} in the code. The fact that they must be
tangent to $M \subset \R^n$ is then only implicitly assumed.  Alternatively,
one can represent a tangent vector at $p \in M$ in the orthonormal basis of $T_p M$
given by the \cpp{orthonormalFrame()} method.


\subsection{Variational problems}

So far, all problems studied using \dunemodule{dune-gfe} have been variational.
By this we mean that the problem takes the form of an energy functional
\begin{equation*}
 J : H(\Omega;M) \to \R,
\end{equation*}
where $H(\Omega;M)$ is a suitable space of functions mapping $\Omega$ to $M$.
Solutions $u$ of the problem are local minimizers of this functional.

Typically, such variational GFE problems are described by an energy density
\begin{equation*}
 f : \Omega \times TM \to \R,
\end{equation*}
and the energy takes the form
\begin{equation}
\label{eq:integral_energy}
 J(v)
 \colonequals
 \int_\Omega f(x, v(x), \nabla v(x))\,dx.
\end{equation}
(We disregard the case of densities depending on higher derivatives.)


\section{Efficient assembly of the tangent problems}

To minimize the functional $J$ in a GFE space requires first and second derivatives
of the energy $J$.%
%
\footnote{Gradient-descent methods would require first derivatives only, but \dunemodule{dune-gfe}
is geared towards second-order methods.}
%
The second derivative will be a sparse symmetric matrix
(the \emph{Hesse} or \emph{tangent} matrix).  As in the case of classical
finite elements, it can be assembled by considering the integral~\eqref{eq:integral_energy}
element by element, and appropriately summing up the element-wise contributions.

We now consider a single element $T$ of a grid of $\Omega$.
A GFE discretization approximates $u$ on $T$ by a function $u_h : T \to M$
that is constructed from interpolating a set of $m$ values $u_1,\dots,u_m \in M$
to a function on $T$. In the following we write $u_h(u_1,\dots,u_m,\cdot)$ to denote
the function obtained by interpolating the coefficients $u_1,\dots,u_m$, and
$u_h(u_1,\dots,u_m,x)$ for the value of that function at $x \in T$. Similarly,
we will write $\nabla u_h$ for the derivative of the interpolation function
with respect to $x$. We can therefore consider the energy on $T$ as a
function on $M^m$:
\begin{align*}
 E & \; : M^m \to \R
 \\
 E(u_1,\dots,u_m)
 & \colonequals
 \int_T f\big(x, u_h(u_1,\dots,u_m,x), \nabla u_h(u_1,\dots,u_m,x) \big)\,dx.
\end{align*}
As usual the integral is approximated by a quadrature rule.  We therefore replace it
by the sum
\begin{align}
\label{eq:integral_energy_quadrature}
 E(u_1,\dots,u_m)
 \approx
 \sum_{(\omega,x) \in Q} f\big(x, u_h(u_1,\dots,u_m,x), \nabla u_h(u_1,\dots,u_m,x) \big),
\end{align}
where $Q$ is a suitable quadrature rule with weights $\omega_1,\dots,\omega_{\abs{Q}} \in \R$
and points $x_1,\dots,x_{\abs{Q}} \in T$.

\subsection{Euclidean vs.\ Riemannian derivatives}

For Newton-type methods one needs the gradient and Hessian of the functional $E$.
However, as the energy $E$ is defined on a Riemannian manifold (the product manifold
$M^m$), we need the Riemannian derivatives instead of the standard Euclidean ones
\cite{absil}.

Conceptually, computing the first and second Riemannian derivatives consists of two steps.
Let first
\begin{equation*}
 \bar{E} : (\R^n)^m \to \R
\end{equation*}
be any smooth extension of $E$ from the manifold $M^m$ into the surrounding space $(\R^m)^n$,
and let $\bar{\nabla} \bar{E}$ and $\bar{\nabla}^2 \bar{E}$ be the standard Euclidean
gradient and Hesse matrix of $E$, respectively. It is well-known (see, e.g., \cite{})
that the Riemannian gradient can be obtained from the Euclidean one by an orthogonal
projection of $\bar{\nabla}\bar{E}$ onto the tangent space of $M^m$
\begin{equation}
\label{eq:gradient_riemannian_from_euclidean}
 \nabla E(u_1,\dots,u_m)
 =
 \mathcal{P}_{(u_1,\dots,u_m)} \bar{\nabla}\bar{E}(u_1,\dots,u_m),
\end{equation}
where for any $(u_1,\dots,u_m) \in M^m$, $\mathcal{P}_{(u_1,\dots,u_m)}$ shall denote
the orthogonal projection from $T_{(u_1,\dots,u_m)}(\R^m)^n$ onto $T_{(u_1,\dots,u_m)} M^m$.
Since $M^m$ is a product manifold, the projection operator $\mathcal{P}_{(u_1,\dots,u_m)}$
is actually a separate projection for each factor space $\mathcal{P}_{u_1}, \dots, \mathcal{P}_{u_m}$.

In~\cite{absil_mahony_trumpf:2013}, \citeauthor{absil_mahony_trumpf:2013} showed a
corresponding formula for computing the Riemannian Hessian from the Euclidean one
\begin{align}
\label{eq:hessian_riemannian_from_euclidean}
 \operatorname{Hess} E(u_1,\dots,u_m)[z]
 & =
 \mathcal{P}_{(u_1,\dots,u_m)} \bar{\nabla}^2 \bar{E}(u_1,\dots,u_m) z
 \\
 \nonumber
 & \quad
 + \mathfrak{A}_{(u_1,\dots,u_m)}(z, \mathcal{P}^\perp_{(u_1,\dots,u_m)}).
\end{align}
Here $\mathfrak{A}_{(u_1,\dots,u_m)}$ is the Weingarten map of $M^m$ at $(u_1,\dots,u_m)$,
which encodes information about the curvature of embedded submanifolds.
\begin{definition}[Weingarten map]
 The Weingarten map of the submanifold $M$ at $x$ is the operator $\mathfrak{A}_x$
 that takes as arguments a tangent vector $z \in T_x M$ and a normal vector $v \in T^\perp_x M$
 and returns the tangent vector $\mathfrak{A}_x (z, v) = -\mathcal{P}_xD_z V$,
 where $V$ is any local extension of $v$ to a normal vector field on $M$.
\end{definition}
\todosander{Definition von $D_z$ fehlt.}
\todosander{Show that this is always a symmetric matrix!}
\todosander{Note: For this construction we need the \emph{embedded} gradient!}


\subsection{Direct assembly of the Euclidean gradient and Hessian}
\label{sec:direct_assembly}

We first show how to compute the Euclidean gradient and Hesse matrix.
The \dunemodule{dune-gfe} code always computes the two together, because
computing them is expensive, and computing the second derivative almost
always produces the first one as a by-product.
The easiest way is to use ADOL-C to compute the tangent problems by feeding
the entire quadrature loop~\eqref{eq:integral_energy_quadrature} as a
single function, and have it
compute the first and second derivatives using the \cpp{jacobian} and
\cpp{hessian2} methods.%
%
\footnote{Always use ADOL-C's vector mode for the second derivatives.
 It is much faster than scalar mode.}
%
This is implemented in the class \cpp{LocalGeodesicFEADOLCStiffness}.
(along with the modifications of Chapter~\ref{sec:euclidean_to_riemannian} that turn
the Euclidean derivatives into Riemannian ones.)
It works, and has been used successfully for all my publications up and
including \cite{nebel_sander_birsan_neff:2023}.

\subsection{Computing Euclidean derivatives by separating the interpolation from the density}

The approach of Chapter~\ref{sec:direct_assembly} is simple but slow.
A part of the slowness is intrinsic: Due to the nonlinear nature of the
GFE interpolation, and the fact that in many cases it cannot be given explicitly,
deriving energies that involve such interpolations is an expensive task.
Run-time can be saved, however, by differentiating the density $f$ and the interpolation function $u_h$
separately. This is what the \cpp{LocalIntegralStiffness} class implements.
Let
\begin{equation*}
 \bar{\nabla} \bar{E}(u_1,\dots,u_m) \in \R^{mn}
\end{equation*}
be the Euclidean gradient of $\bar{E}$ at $u_1, \dots, u_m \in \R^n$, and
\begin{equation*}
 \bar{\nabla}^2 \bar{E}(u_1,\dots,u_m) \in \R^{mn \times mn}
\end{equation*}
the corresponding Hesse matrix.  With the chain rule we get
\begin{equation*}
 \bar{\nabla} \bar{E}(u_1,\dots,u_m)
 =
 \sum_{(x,\omega) \in Q} \mathcal{J}^T_{u_h,x} \nabla f(x, u_h(x), \nabla u_h(x)).
\end{equation*}
Here,
\begin{equation*}
 \mathcal J_{u_h,x}
 \colonequals
 \begin{pmatrix}
    \frac{\partial u_h(x)}{\partial u_1}, \dots \frac{\partial u_h(x)}{\partial u_m} \\
    \frac{\partial \nabla u_h(x)}{\partial u_1},
        \dots, \frac{\partial \nabla u_h(x)}{\partial u_m}
 \end{pmatrix}
 \in
 \R^{(n+dn) \times mn}
\end{equation*}
is the Jacobi matrix of the GFE interpolation $(u_1,\dots, u_m) \mapsto (u_h(x), \nabla u_h(x))$
for fixed $x \in T$, and
\begin{equation*}
 \nabla f(x, u_h(x), \nabla u_h(x))
 \in
 \R^{n+dn}
\end{equation*}
is the derivative of the density with respect to the second and third argument.

To understand the sizes of these objects, note that a pair $(u_h(x),\nabla u_h(x))$
can be interpreted as a vector in $\R^n \times \R^{dn} \simeq \R^{n+dn}$
(with $d = \dim \Omega$), and the tuple of coefficients $u_1,\dots,u_m$
can be interpreted as a vector in $\R^{mn}$.

Likewise, for the second derivative of $E$ we get
\begin{align*}
 \bar{\nabla}^2 \bar{E}(u_1,\dots,u_m)
 & =
 \sum_{(x,\omega) \in Q} \mathcal J_{u_h,x}^T \nabla^2 f(x,u_h(x),\nabla u_h(x)) \mathcal{J}_{u_h,x}
 \\
 %
 & \qquad \qquad +
 \sum_{(x,\omega) \in Q} \nabla f(x,u_h(x),\nabla u_h(x)) \mathcal H_{u_h,x}.
\end{align*}
where
\begin{align*}
 (\mathcal{H}_{u_h,x})_{ijk}
 \colonequals
 \frac{\partial^2 (u_h(x))_i}{\partial (u_{j/n})_{j \operatorname{mod} n} \partial (u_{k/n})_{k \operatorname{mod} n}}
 \qquad
 i = 1,\dots,n
 \qquad
 j,k = 1, \dots, mn
 \shortintertext{and}
 (\mathcal{H}_{u_h,x})_{ijk}
 \colonequals
 \frac{\partial^2 (\nabla u_h(x))_i}{\partial (u_{j/n})_{j \operatorname{mod} n} \partial (u_{k/n})_{k \operatorname{mod} n}}
 \qquad
 i = n+1,\dots,n+dn
 \quad
 j,k = 1, \dots, mn
\end{align*}
The product of the vector $\nabla f$ and the three-indices-object $\mathcal H_{u_h,x}$
is to be interpreted as a contraction with respect to the first index of $\mathcal H_{u_h,x}$:
\begin{equation*}
 (\nabla f \mathcal H_{u_h,x})_{jk}
 \colonequals
 \sum_{i=1}^{n+dn}
 (\nabla f)_i (\mathcal H_{u_h,x})_{ijk}
 \in
 \R^{mn \times mn}.
\end{equation*}

In this representation, the derivatives of $f$ and of $u_h$ are separated and
can be computed separately. When done right this saves quite a bit of run-time.

\subsection{From Euclidean to Riemannian derivatives}
\label{sec:euclidean_to_riemannian}

In the second step the Euclidean gradient and Hessian are modified according to
\eqref{eq:gradient_riemannian_from_euclidean} and~\eqref{eq:hessian_riemannian_from_euclidean},
respectively, to obtain the Riemannian counterparts.
Calling this a second step is a conceptual distinction, because the code
does not actually compute the Euclidean Hessian before modifying it.
\todosander{Explain why!}
For the gradient, this modification is simply the multiplication from the left
with the projection operator $\mathcal{P}$. In practice, two things happen:
Applying the projection operator to a vector in $T_x M \simeq \R^n$ produces a
vector that is tangent to $M$, but still represented as a vector in $\R^n$
(in \dunemodule{dune-gfe} parlance: an embedded tangent vector,
see Chapter~\ref{sec:representation_of_manifolds}).

\todosander{Den Rest schreiben!}

\subsection{Computing the derivatives $\mathcal{H}$ and $\mathcal{H}$ of the geometric interpolation}

\begin{remark}
 $\mathcal H$ is an object with three indices, which is potentially memory- and
 time-consuming to compute and store.  Note, however, that what is needed is not
 $\mathcal H$ itself but its contraction with the vector $\nabla f(x,u_h(x),\nabla u_h(x))$.
 AD systems that implement reverse-mode AD (such as ADOL-C) can compute this
 contracted second derivative directly (in fact, such contractions is all they can compute).
 The vector to contract with is called \emph{adjoint} in AD parlance.
 (Specifically, \dunemodule{dune-gfe} uses the \texttt{hos\_ov\_reverse} method).
 No full third-order object is ever computed.
\end{remark}

\begin{remark}
 If the target manifold $M$ happens to be the Euclidean space, then
 \begin{equation*}
  u_h(x) = \sum_i u_i \varphi_i(x)
 \end{equation*}
 for a set of shape functions $\varphi_i$.  While Euclidean spaces by themselves
 are not interesting in a GFE context, they do appear frequently as factors in
 larger spaces (e.g., Cosserat problems typically use $M = \R^3 \times SO(3)$.

 \todosander{Hier nur den skalaren Fall erklären?  Besser nicht...}
 Consequently, in that case
 \begin{align*}
  \mathcal{J} = \Big( \varphi_1(x), \dots, \varphi_m(x), \nabla \varphi_1(x), \dots \nabla \varphi_m(x)\Big)
 \end{align*}
 and
 \begin{equation*}
  \mathcal{H} = 0.
 \end{equation*}
 The method that computes $\mathcal J$ and $\mathcal H$ is therefore specialized
 for this case, and hand-codes these expressions.  In my experience, the time gained
 by not using AD for this case is rather small, though.
\end{remark}

\begin{remark}
 If $M$ is a product space, then $\mathcal J$ and $\mathcal H$ are block diagonal.
 This must be exploited by the implementation.
\end{remark}

\begin{remark}
 Besides the trivial case $M = \R^n$, the computation of $\mathcal J$ and $\mathcal H$
 can also be hand-implemented for other cases.  One example is projection-based
 interpolation for the spheres $M = S^{n-1}$.
\end{remark}


\section{The derivatives of projection-based interpolation in the sphere}

For the projection-based interpolation mapping into the unit sphere, the derivatives
required by the geometric finite element method can be computed by hand
with reasonable effort. The \dunemodule{dune-gfe} module contains these hand-coded
derivatives because they are faster than computing the derivatives with ADOL-C.

We document here the expressions for the derivatives. They are implemented in the file
\texttt{interpolationderivatives.hh}, in a specialization of the class
\cpp{InterpolationDerivatives}.

Let $\varphi_i : T \to \R$ be a set of weight functions that sum up to~$1$
everywhere, and let $v_1,\dots,v_m \in S^{n-1} \subset \R^n$ be values
on the unit sphere.  Recall that the projection-based interpolation at
a point $\xi \in T$ is
\begin{equation*}
 I(v_1,\dots,v_m;\xi)
 \colonequals
 P(\omega(\xi)),
\end{equation*}
where
\begin{equation*}
 \omega(\xi) \colonequals \sum_{i=1}^m \varphi_i(\xi)v_i
\end{equation*}
is the Euclidean interpolation between the values $v_1,\dots,v_m$, and
\begin{equation*}
 P : \R^n \to S^{n-1}
 \qquad
 P(x) \colonequals \frac{x}{\abs{x}}
\end{equation*}
is the radial projection onto the unit sphere.

\begin{remark}
 While this is the definition of projection-based interpolation as given in articles
 such as~\cite{grohs_hardering_sander_sprecher:2019} and \cite{hardering_sander:2020},
 the code currently implements something subtly different.  Specifically, the
 copy constructor and assignment operators of \texttt{UnitVector} project their
 arguments to $S^{n-1}$, which changes the definition of $\omega(\xi)$ given above to
 \begin{equation*}
  \omega(\xi)
  \colonequals
  \sum_{i=1}^m \varphi_i(\xi) \frac{v_i}{\abs{v_i}}.
 \end{equation*}
 While this does not change the value of the projection-based interpolation as long as
 the input values $v_1,\dots,v_m$ really do have unit length, the chain rule makes
 additional terms appear in the derivative.

 I am not actually convinced that having the copy constructor and assignment perform
 a projection is a good idea, and I may change this in the future.
\end{remark}


\subsection{Derivatives of the projection onto the sphere}

By the chain rule, derivatives of $I$ involve derivatives of the projection operator~$P$.
We therefore start by computing these derivatives.

Before starting, note that for any $N \in \mathbb{Z}$ and $j=1,\dots,n$
\begin{equation*}
 \frac{\partial}{\partial x_j} \abs{x}^N
 =
 \frac{\partial}{\partial x_j} \Big( \sum_k x_k^2 \Big)^{\frac{N}{2}}
 =
 \frac{N}{2} \abs{x}^{N-2} \sum_k \frac{\partial}{\partial x_j} x_k^2
 =
 N x_j \abs{x}^{N-2}.
\end{equation*}

\subsubsection{Partial derivatives}

The partial derivatives do not actually appear in the implementation, but they are
an instructive intermediate step.

\paragraph{First derivatives}

The Jacobian of $P$ is the matrix
 \begin{equation*}
  (\nabla P)_{ij}
  =
  \frac{\partial}{\partial x_j} \frac{x_i}{\abs{x}}
  =
  \frac{\frac{\partial x_i}{\partial x_j}\abs{x} - x_i \frac{\partial}{\partial x_j} \abs{x}}{\abs{x}^2}
  =
  \delta_{ij} \frac{1}{\abs{x}} - \frac{x_i x_j}{\abs{x}^3}.
 \end{equation*}

\paragraph{Second derivatives}
The second derivative is the third-order tensor
\begin{align*}
 (\nabla^2 P)_{ijk}
 & =
 \frac{\partial}{\partial x_j \partial x_k} \frac{x_i}{\abs{x}}
 \\
 & =
 \frac{\partial}{\partial x_k} \Big[ \frac{\delta_{ij}}{\abs{x}} - \frac{x_i x_j}{\abs{x}^3} \Big]
 \\
 & =
 \delta_{ij} \frac{\partial}{\partial x_k} \frac{1}{\abs{x}}
   - \frac{\partial}{\partial x_k}\Big( x_i x_j \frac{1}{\abs{x}^3} \Big)
 \\
 & =
 -\delta_{ij} \frac{x_k}{\abs{x}^3} - \delta_{ik} x_j \frac{1}{\abs{x}^3}
   - \delta_{jk} x_i \frac{1}{\abs{x}^3} + 3 x_i x_j \frac{x_k}{\abs{x}^5}
 \\
 & =
 - \Big( \delta_{ij}x_k + \delta_{ik}x_j + \delta_{jk}x_i\Big) \frac{1}{\abs{x}^3}
 + 3 \frac{x_i x_j x_k}{\abs{x}^5}.
\end{align*}


\paragraph{Third derivatives}

The third derivative is a fourth-order object
\begin{align*}
 (\nabla^3 P)_{ijkl}
 & =
 \frac{\partial}{\partial x_l}
 \bigg[ - \Big( \delta_{ij}x_k + \delta_{ik}x_j + \delta_{jk}x_i\Big) \frac{1}{\abs{x}^3}
 + 3 \frac{x_i x_j x_k}{\abs{x}^5} \bigg]
 \\
 & =
 (-\delta_{ij}\delta_{kl} - \delta_{ik}\delta_{jl} - \delta_{jk} \delta_{il}) \frac{1}{\abs{x}^3}
 + (-\delta_{ij}x_k - \delta_{ik}x_j - \delta_{jk}x_i)(-3)\frac{x_l}{\abs{x}^5}
 \\
 & \qquad
 + 3\Big(\frac{\delta_{il} x_j x_k}{\abs{x}^5} + \frac{\delta_{jl} x_i x_k}{\abs{x}^5}
   + \frac{\delta_{kl} x_i x_j}{\abs{x}^5} + x_i x_j x_k (-5) \frac{x_l}{\abs{x}^7} \Big)
 \\
 & =
 (\delta_{ij}\delta_{kl} + \delta_{ik}\delta_{jl} + \delta_{jk} \delta_{il}) \frac{-1}{\abs{x}^3}
 \\
 & \qquad
 + \frac{3}{\abs{x}^5} \Big[ \delta_{ij}x_k x_l + \delta_{ik}x_jx_l + \delta_{jk}x_ix_l
   + \delta_{il} x_j x_k + \delta_{jl} x_i x_k + \delta_{kl}x_i x_j \Big]
 \\
 & \qquad
 + x_i x_j x_k x_l \frac{-15}{\abs{x}^7}.
\end{align*}

\subsubsection{Directional derivatives and contractions}

As can be seen below, the final derivative formulas use the derivatives of $P$
by contracting it with different vectors.  These contractions can be computed
relatively cheaply, because the derivatives of $P$ consist only of diagonal
and low-rank terms.

\paragraph{First derivative}

Let $w \in \R^n$ be an adjoint vector and $\delta^1 \in \R^n$ a direction.  Then
(with Einstein summation)
\begin{equation*}
 \nabla P[w,\delta^1]
 =
 (\nabla P)_{ij} w_i \delta^1_j
 =
 \frac{\langle w,\delta^1\rangle}{\abs{x}} - \frac{\langle w,x\rangle \langle \delta^1,x\rangle}{\abs{x}^3}.
\end{equation*}

\paragraph{Second derivative}

Let $w \in \R^n$ be an adjoint vector, and $\delta^1, \delta^2 \in \R^n$ directions.
Then
\begin{align*}
 \nabla^2P[w,\delta^1,\delta^2]
 & =
 (\nabla^2 P)_{ijk} w_i \delta^1_j \delta^2_k
 \\
 & =
 \big(\langle w, \delta^1 \rangle \langle x,\delta^2 \rangle
    + \langle w, \delta^2 \rangle \langle x,\delta^1 \rangle
    + \langle \delta^1, \delta^2\rangle \langle w,x \rangle \big)\frac{-1}{\abs{x}}
 + 3\langle w,x \rangle \langle \delta^1,x \rangle \langle \delta^2,x \rangle \frac{3}{\abs{x}^5}.
\end{align*}


TODO: WIRD DIE FOLGENDE GEMISCHTE VARIANTE EVENTUELL GEBRAUCHT?
\begin{align*}
 (w \nabla^2 P)_{jk}
 & =
 \sum_{i=1}^n w_i (\nabla^2 P)_{ijk}
 \\
 & =
 w_i\bigg[- \Big( \delta_{ij}x_k + \delta_{ik}x_j + \delta_{jk}x_i\Big) \frac{1}{\abs{x}^3}
 + 3 \frac{x_i x_j x_k}{\abs{x}^5} \bigg]
 \\
 & =
 \big( - w_j x_k - w_kx_j - \langle w,x\rangle \delta_{jk}\big)\frac{1}{\abs{x}}
 + 3\langle w,x \rangle \frac{x_j x_k}{\abs{x}^5}.
\end{align*}

\paragraph{Third derivative}

Let $w \in \R^n$ be an adjoint vector, and $\delta^1, \delta^2, \delta^3 \in \R^n$ directions.
Then
\begin{align*}
 \nabla^3P[w,\delta^1,\delta^2,\delta^3]
 & =
 (\nabla^3 P)_{ijkl} w_i \delta^1_j \delta^2_k \delta^3_l
 \\
 & =
 \big(\langle w, \delta^1 \rangle \langle \delta^2,\delta^3 \rangle
    + \langle w, \delta^2 \rangle \langle \delta^1,\delta^3 \rangle
    + \langle w, \delta^3 \rangle \langle \delta^1,\delta^2 \rangle \big)
    \frac{-1}{\abs{x}^3}
 \\
 & \qquad
 \Big( \langle w,\delta^1 \rangle \langle x, \delta^2 \rangle \langle x, \delta^3 \rangle
  + \langle w,\delta^2 \rangle \langle w, \delta^2 \rangle \langle w, \delta^3 \rangle
  + \langle \delta^1,\delta^2 \rangle \langle x, w \rangle \langle x, \delta^3 \rangle
  \\
 & \qquad
  + \langle w,\delta^3 \rangle \langle x, \delta^2 \rangle \langle x, \delta^1 \rangle
  + \langle \delta^1,\delta^3 \rangle \langle w, x \rangle \langle x, \delta^2 \rangle
  + \langle \delta^2,\delta^3 \rangle \langle w, x \rangle \langle x, \delta^1 \rangle \Big)
 \frac{3}{\abs{x}^5}
 \\
  & \qquad
 \langle w,x \rangle \langle \delta^1,x \rangle \langle \delta^2,x \rangle \langle \delta^3,x \rangle
 \frac{-15}{\abs{x}^7}.
\end{align*}

Actually we know that at least some of the directions are orthogonal, but this is not
currently used by the code.

TODO: WIRD DIE FOLGENDE GEMISCHTE VARIANTE EVENTUELL GEBRAUCHT?
The contraction with a vector $w \in \R^n$ is
\begin{align*}
 w_i (\nabla^3 P)_{ijkl}
 & =
 \big( w_j \delta_{kl} + w_k \delta_{jl} + w_l \delta_{jk} \big) \frac{-1}{\abs{x}^3}
 \\
 & \quad
 + \frac{3}{\abs{x}^5}
 \Big[ w_j x_k x_l + w_k x_j x_l + \langle w,x \rangle x_l \delta_{jk} + w_l x_j x_k
    + \langle w,x \rangle \delta_{jl} x_k + \langle w,x \rangle \delta_{kl} x_j \Big]
 \\
 & \quad
 - \frac{15}{\abs{x}^7} \langle w,x \rangle x_j x_k x_l.
\end{align*}



\subsection{Derivatives of the interpolation value}

We now compute the derivatives of
\begin{equation*}
 I(v_1,\dots,v_m; \xi)
 =
 P(\omega(\xi))
 =
 P\Big(\sum_{i=1}^m v_i \varphi_i(\xi)\Big)
\end{equation*}
with respect to the coefficients $v_j$, for fixed $\xi$.

\subsubsection{Partial derivatives}

We begin with partial derivatives.  Interpret the $k$-th coefficient $v_k$
as a vector in $\R^n$.

Then the partial derivative is a third-order object: We compute the first derivative
of the $\alpha$-th component of the vector $P(\omega) \in S^{n-1} \subset \R^n$
with respect to the $l$-th component of the $k$-th data point $v_k$
\begin{align*}
 \frac{\partial}{\partial v_k^l} \big( P (v) \big)_\alpha
 & =
 \frac{\partial P(v)_\alpha}{\partial v^m} \cdot \frac{\partial v^m}{\partial v_k^l}
 \\
 & =
 (\nabla P(v))_{\alpha m} \cdot \frac{\partial v^m}{\partial v_k^l}
 \\
 & =
 (\nabla P(v))_{\alpha m} \cdot \delta_{ml} \varphi_k
 =
 (\nabla P(v))_{\alpha l} \varphi_k.
\end{align*}

The code numbers all coefficient components $v_{1,1},\dots,v_{m,n}$ with a single
flat index. Therefore, as a data structure the derivative is a matrix
with $n$ rows and $mn$ columns.

\bigskip

The second derivative is
\begin{equation*}
 \frac{\partial^2}{\partial v_j \partial v_k} I(v_1,\dots,v_m;\xi)
 =
 \frac{\partial}{\partial v_k} \nabla P(v) \varphi_j (\xi)
 =
 \varphi_k \nabla^2 P(v) \varphi_j.
\end{equation*}
In index notation, this is an object with five indices
\begin{equation*}
 \frac{\partial}{\partial v_\gamma^\delta} \frac{\partial}{\partial v_k^l} \big(P(v)\big)_\alpha
 =
 \frac{\partial}{\partial v_\gamma^\delta} \Big[ (\nabla P)_{\alpha l} \varphi_k \Big]
 =
 (\nabla^2 P)_{\alpha l \delta} \varphi_k \varphi_\gamma.
\end{equation*}



\subsection{Derivatives of the derivative with respect to $\xi$}

\begin{equation*}
 \nabla I(v;\xi)
 =
 \frac{\partial}{\partial \xi} I(v;\xi)
 =
 \nabla P(v) \sum_{i=1}^m v_i \nabla \varphi_i
 =
 \nabla P(v) \nabla v.
\end{equation*}
This is a matrix with entries
\begin{align*}
 (\nabla I)_{\alpha \beta}
 & =
 \frac{\partial \big(P(v)\big)_\alpha}{\partial \xi_\beta}
 \\
 & =
 \frac{\partial P(v)_\alpha}{\partial v^m} \cdot \frac{\partial v^m}{\partial \xi_\beta}
 \\
 & =
 \frac{\partial P(v)_\alpha}{\partial v^m} \cdot \frac{\partial \sum_i v_i^m \varphi_i(\xi)}{\partial \xi_\beta}
 \\
 & =
 \frac{\partial P(v)_\alpha}{\partial v^m} \cdot \sum_i v_i^m \frac{\partial \varphi_i(\xi)}{\partial \xi_\beta}
\end{align*}



The first derivative of this with respect to a coefficients is
\begin{equation*}
 \frac{\partial^2}{\partial \xi \partial v_j} I(v_1,\dots,v_m;\xi)
 =
 \frac{\partial}{\partial v_j} \Big[ \nabla P(v) \sum_{i=1}^m v_i \nabla \varphi_i \Big]
 =
 \varphi_j \nabla^2P(v) \sum_{i=1}^m v_i \nabla \varphi_i + \nabla P(v) \nabla \varphi_j
\end{equation*}

In coefficients:
\begin{align*}
 \frac{\partial^2 I}{\partial \xi_\beta \partial_i^j}
 & =
 \frac{\partial}{\partial v_i^j} \Big[ (\nabla P(v))_{\alpha m} \cdot (\nabla v^m)_\beta \Big]
 \\
 & =
 \Big( \frac{\partial}{\partial v_i^j} (\nabla P(v))_{\alpha m} \Big) \cdot (\nabla v^m)_\beta
 +
 (\nabla P(v))_{\alpha m} \cdot \Big( \frac{\partial}{\partial v_i^j} (\nabla v^m)_\beta \Big)
 \\
 & =
 \big( \nabla^2 P(v) \big)_{\alpha m j} \cdot \varphi_i \cdot ( \nabla v^m)_\beta
 +
 (\nabla P(v))_{\alpha m} \cdot \delta_{jm} (\nabla \varphi_i)\beta
 \\
 & =
 \big( \nabla^2 P(v) \big)_{\alpha m j} \cdot \varphi_i \cdot ( \nabla v^m)_\beta
 +
 (\nabla P(v))_{\alpha j} \cdot (\nabla \varphi_i)\beta.
\end{align*}


The second derivative is
\begin{align*}
 \frac{\partial^3}{\partial \xi \partial v_j \partial v_k} I(v_1,\dots,v_m;\xi)
 & =
 \frac{\partial}{\partial v_k} \Big[ \varphi_j \nabla^2P(v) \sum_{i=1}^m v_i \nabla \varphi_i + \nabla P(v) \nabla \varphi_j \Big] \\
 & =
 \varphi_j \nabla^3 P(v) \varphi_k \nabla v + \varphi_k \nabla^2 P(v) \nabla \varphi_k + \nabla^2P(v) \varphi_k \nabla \varphi_j.
\end{align*}

In coefficients:

\begin{align*}
 \frac{\partial^3 I_\alpha}{\partial \xi_\beta \partial v_i^j \partial v_k^l}
 & =
 \frac{\partial}{\partial v_k^l} \Big[ \big(\nabla^2 P(v)\big)_{amj} \varphi_i \Big] \cdot (\nabla v^m)_\beta
 \\
 & \qquad
 + \big(\nabla^2 P(v)\big)_{\alpha m j} \varphi_i \frac{\partial}{\partial v_k^l} (\nabla v^m)_\beta
 \\
 & \qquad
 + \frac{\partial}{\partial v_k^l} \Big[ \big(\nabla P(v)\big)_{\alpha m} \Big]
  \cdot \delta_{jm} \cdot (\nabla \varphi_i)_\beta
 \\
 & =
 \big(\nabla^3 P(v) \big)_{\alpha m j l} \varphi_i \varphi_k \cdot \frac{\partial v^m}{\partial \xi_\beta}
 \\
 & \qquad
 + \big(\nabla^2 P(v) \big)_{\alpha m j} \varphi_i \cdot \delta_{ml} \frac{\partial \varphi_k}{\partial \xi_\beta}
 + \big(\nabla^2 P(v) \big)_{\alpha m l} \varphi_k \cdot \delta_{mj} \frac{\partial \varphi_i}{\partial \xi_\beta}
 \\
 & =
 \big(\nabla^3 P(v) \big)_{\alpha m j l} \varphi_i \varphi_k \cdot \frac{\partial v^m}{\partial \xi_\beta}
 + \big(\nabla^2 P(v) \big)_{\alpha j l}
 \Big [\varphi_i \cdot \frac{\partial \varphi_k}{\partial \xi_\beta}
 + \varphi_k \cdot \frac{\partial \varphi_i}{\partial \xi_\beta} \Big]
\end{align*}




\printbibliography

\end{document}
